\subsection{The scale of low-credibility AI disclosure}

Figure~\ref{fig:mismatch_surge} documents the central descriptive pattern. Low-credibility AI disclosure rises sharply after 2022. The number of AI-talking firm-years flagged as PatentMismatch rises sharply after that date, and the mismatch share among AI-talking firm-years reaches 41.48\% in 2024 and 42.12\% in 2025. The expansion of AI language is not limited to disclosure volume. It also reflects a rise in firm-years where disclosure composition is inconsistent with contemporaneous AI patenting.

Table~\ref{tab:summary_stats} describes the annual ever-speaker panel that underlies the main tests. The panel contains 50,840 firm-year observations across 5,084 firms, of which 13,777 are AI-talking firm-years, and the classified corpus contains 147,879 AI-related sentences. AI-related disclosure is highly skewed, while AI patenting remains concentrated in a much smaller subset of firms. The surge in AI-related disclosure occurs while visible implementation remains uneven.

\begin{table}[!htbp]
\centering
\scriptsize
\renewcommand{\arraystretch}{0.92}
\caption{Expanded-panel summary statistics and sample coverage}
\label{tab:summary_stats}

\begin{minipage}{\textwidth}
\textit{Panel A. Coverage and corpus totals}
\end{minipage}
\vspace{0.2em}

\begin{tabular}{>{\raggedright\arraybackslash}p{0.28\textwidth}>{\raggedright\arraybackslash}p{0.12\textwidth}>{\raggedright\arraybackslash}p{0.48\textwidth}}
\hline
Metric & Value & Notes \\
\hline
Ever-speaker firms & 5,084 & Unique CIKs in the canonical annual panel. \\
Firm-year observations & 50,840 & Annual panel coverage for 2016--2025. \\
AI-talking firm-years & 13,777 & Firm-years with at least one classified AI sentence. \\
Total classified AI sentences & 147,879 & Corpus total across the 2016--2025 ever-speaker panel. \\
Actionable AI sentences & 67,080 & Corpus total. \\
Speculative AI sentences & 17,900 & Corpus total. \\
Irrelevant AI sentences & 62,899 & Corpus total. \\
Firm-years with Compustat controls & 31,291 & Non-missing log assets in the annual panel. \\
Firm-years with matched market cap & 23,747 & Annual rows with merged CRSP market-cap data. \\
Filing-event observations & 7,355 & Distinct filing events across 2,753 firms. \\
Filing events with CAR[$-1,+1$] & 7,342 & Usable immediate-return events. \\
Filing events with BHAR[$+2,+63$] & 7,098 & Usable 3-month drift events. \\
Filing events with BHAR[$+2,+252$] & 4,756 & Usable 12-month drift events. \\
\hline
\end{tabular}

\vspace{0.5em}
\begin{minipage}{\textwidth}
\textit{Panel B. Annual-panel variables}
\end{minipage}
\vspace{0.2em}

\resizebox{\textwidth}{!}{%
\begin{tabular}{lcccccccc}
\hline
Variable & Mean & Std. Dev. & p5 & p25 & p50 & p75 & p95 & N \\
\hline
Total AI sentences per firm-year & 2.909 & 17.306 & 0.000 & 0.000 & 0.000 & 1.000 & 13.000 & 50,840 \\
Actionable AI sentences per firm-year & 1.319 & 8.293 & 0.000 & 0.000 & 0.000 & 0.000 & 6.000 & 50,840 \\
Speculative AI sentences per firm-year & 0.352 & 2.494 & 0.000 & 0.000 & 0.000 & 0.000 & 1.000 & 50,840 \\
Irrelevant AI sentences per firm-year & 1.237 & 8.209 & 0.000 & 0.000 & 0.000 & 0.000 & 6.000 & 50,840 \\
AI Focus & 0.463 & 0.933 & 0.000 & 0.000 & 0.000 & 0.693 & 2.639 & 50,840 \\
Actionable share & 0.129 & 0.305 & 0.000 & 0.000 & 0.000 & 0.000 & 1.000 & 50,840 \\
SpecShare & 0.031 & 0.136 & 0.000 & 0.000 & 0.000 & 0.000 & 0.200 & 50,840 \\
CredAI & 0.000 & 0.755 & -0.298 & -0.018 & -0.018 & -0.018 & 0.344 & 50,840 \\
Actionable-to-speculative ratio ($A_S$) & 0.210 & 0.516 & 0.000 & 0.000 & 0.000 & 0.000 & 1.386 & 50,840 \\
PatentMismatch & 0.096 & 0.295 & 0.000 & 0.000 & 0.000 & 0.000 & 1.000 & 50,840 \\
AI patents & 0.160 & 2.969 & 0.000 & 0.000 & 0.000 & 0.000 & 0.000 & 50,840 \\
Total patents & 6.012 & 76.473 & 0.000 & 0.000 & 0.000 & 0.000 & 10.000 & 50,840 \\
Log assets & 6.981 & 2.799 & 2.096 & 5.361 & 7.317 & 8.889 & 10.992 & 31,291 \\
Leverage & 0.241 & 0.253 & 0.000 & 0.026 & 0.185 & 0.371 & 0.708 & 31,186 \\
Cash/assets & 0.238 & 0.275 & 0.004 & 0.037 & 0.114 & 0.352 & 0.879 & 31,290 \\
R\&D/assets & 0.164 & 0.309 & 0.000 & 0.010 & 0.060 & 0.187 & 0.639 & 17,934 \\
CAPX/assets & 0.028 & 0.039 & 0.000 & 0.002 & 0.014 & 0.038 & 0.104 & 31,113 \\
ROA & -0.232 & 0.912 & -1.199 & -0.147 & 0.010 & 0.049 & 0.150 & 31,228 \\
Sales growth & 1.310 & 31.591 & -0.378 & -0.025 & 0.072 & 0.212 & 1.016 & 28,010 \\
Employees & 13.495 & 59.888 & 0.007 & 0.182 & 1.400 & 7.600 & 57.000 & 29,404 \\
Log MktCap/assets & 6.811 & 1.187 & 4.675 & 6.088 & 6.894 & 7.638 & 8.679 & 23,444 \\
Log Q proxy & 6.809 & 1.186 & 4.675 & 6.087 & 6.890 & 7.635 & 8.673 & 23,369 \\
$\Delta$ log Q $t+1$ & -0.062 & 0.541 & -1.057 & -0.305 & -0.024 & 0.216 & 0.757 & 19,894 \\
$\Delta$ shares $t+1$ & 0.083 & 0.407 & -0.097 & -0.010 & 0.006 & 0.049 & 0.562 & 20,391 \\
Issue $>$ 5\% $t+1$ & 0.247 & 0.431 & 0.000 & 0.000 & 0.000 & 0.000 & 1.000 & 20,391 \\
\hline
\end{tabular}%
}

\vspace{0.5em}
\begin{minipage}{\textwidth}
\textit{Panel C. Event-study and market variables}
\end{minipage}
\vspace{0.2em}

\resizebox{\textwidth}{!}{%
\begin{tabular}{lcccccccc}
\hline
Variable & Mean & Std. Dev. & p5 & p25 & p50 & p75 & p95 & N \\
\hline
Total AI sentences per filing & 7.205 & 17.075 & 1.000 & 1.000 & 2.000 & 7.000 & 28.000 & 7,355 \\
Actionable AI sentences per filing & 3.875 & 9.640 & 0.000 & 0.000 & 1.000 & 4.000 & 16.000 & 7,355 \\
Speculative AI sentences per filing & 0.878 & 2.602 & 0.000 & 0.000 & 0.000 & 1.000 & 4.000 & 7,355 \\
Actionable share per filing & 0.545 & 0.415 & 0.000 & 0.000 & 0.600 & 1.000 & 1.000 & 7,355 \\
Speculative share per filing & 0.123 & 0.254 & 0.000 & 0.000 & 0.000 & 0.125 & 1.000 & 7,355 \\
CAR[$-1,+1$] & 0.000 & 0.116 & -0.159 & -0.037 & -0.001 & 0.033 & 0.159 & 7,342 \\
BHAR[$+2,+21$] & -0.022 & 0.219 & -0.315 & -0.109 & -0.023 & 0.051 & 0.244 & 7,265 \\
BHAR[$+2,+63$] & -0.036 & 0.380 & -0.499 & -0.204 & -0.050 & 0.090 & 0.430 & 7,098 \\
BHAR[$+2,+126$] & -0.064 & 0.530 & -0.695 & -0.304 & -0.084 & 0.111 & 0.572 & 6,829 \\
BHAR[$+2,+252$] & -0.043 & 1.384 & -0.902 & -0.452 & -0.135 & 0.151 & 0.919 & 4,756 \\
\hline
\end{tabular}%
}
\caption*{\scriptsize \textit{Note.} Panel A reports sample coverage and corpus totals. Panels B and C report distributional statistics for annual-panel and filing-event variables. Sentence totals in Panel A are corpus totals; sentence variables in Panels B and C are per-firm-year or per-filing quantities.}
\end{table}


\subsection{Validation against later AI realization}

Validation begins with later AI realization. Table~\ref{tab:construct_variant_screen} compares alternative construct definitions. The canonical grant-based mismatch measure predicts weaker subsequent AI patent grants at $t+1$ with a coefficient of $-0.0435$ ($p<0.01$), whereas a stricter mismatch variant is weaker and unstable. A weak-patent-only indicator is stronger for near-term grants by construction, but it is not a substitute for PatentMismatch because it removes the disclosure side of the construct. Canonical mismatch is associated with stronger AI application activity at $t+2$, consistent with delayed realization or experimentation despite weak near-term patent grants. The construct-variant screen supports the canonical disclosure-plus-implementation measure and clarifies the limits of looser alternatives.

\begin{table}[!htbp]
\centering
\caption{Construct-variant screen for future AI realization}
\label{tab:construct_variant_screen}
\caption*{\footnotesize \textit{Note.} This table compares alternative low-credibility disclosure constructs against future AI realization outcomes. Panel A uses future AI patent grants and Panel B uses future AI patent applications. All outcomes are measured at annual horizons relative to the disclosure year. Strict grant mismatch requires both low AS ratio and high speculative share before applying the patent-weakness condition; application mismatch replaces grant weakness with AI patent-application weakness; low credibility only uses the language gate alone; weak patent only uses the patent-weakness condition without the disclosure gate. The table is intended to show why the canonical grant-based PatentMismatch measure is retained for the main analysis. *, **, and *** denote significance at the 10\%, 5\%, and 1\% levels, respectively.}
\begin{tabular}{lcc}
\hline
\multicolumn{3}{l}{\textit{Panel A. Future AI patent-grant outcomes}} \\
Variant & Log(1 + AI patent grants) t+1 & Log(1 + AI patent grants) t+2 \\
\hline
Canonical grant mismatch & -0.0435*** (0.0134) & 0.0085 (0.0081) \\
Strict grant mismatch & -0.0274 (0.0222) & -0.0301 (0.0187) \\
Application mismatch & -0.0308*** (0.0092) & 0.0062 (0.0085) \\
Low credibility only & -0.0034 (0.0131) & 0.0121 (0.0138) \\
Weak patent only & -0.1276*** (0.0239) & 0.0020 (0.0152) \\
Outcome mean & 0.104 & 0.092 \\
Observations & 10129 & 10129 \\
\hline
\multicolumn{3}{l}{\textit{Panel B. Future AI application outcomes}} \\
Variant & Log(1 + AI applications) t+1 & Log(1 + AI applications) t+2 \\
\hline
Canonical grant mismatch & 0.0221* (0.0114) & 0.0685*** (0.0141) \\
Strict grant mismatch & -0.0130 (0.0149) & -0.0032 (0.0167) \\
Application mismatch & -0.0100 (0.0086) & 0.0157* (0.0089) \\
Low credibility only & 0.0084 (0.0134) & 0.0185 (0.0154) \\
Weak patent only & 0.0521*** (0.0182) & 0.1638*** (0.0232) \\
Outcome mean & 0.140 & 0.118 \\
Observations & 10129 & 10129 \\
\hline
\end{tabular}
\end{table}


\subsection{Real-outcome dynamics}

Operating outcomes provide a second validation check. Table~\ref{tab:real_outcome_dynamics_v31} links PatentMismatch to future ROA, R\&D intensity, and related outcomes. In the AI-talking sample, PatentMismatch predicts weaker profitability at the two-year horizon and stronger later R\&D intensity. The coefficient on ROA at $t+2$ is $-0.0672$ ($p<0.05$), while the coefficient on R\&D-to-assets at $t+2$ is $0.0178$ ($p<0.05$). LowCredibility and ApplicationMismatch show similar patterns. The estimates imply a clear pattern: low-credibility AI disclosure is followed by weaker near-term performance and higher subsequent R\&D.

\begin{table}[!htbp]
\centering
\caption{Real-outcome dynamics in the AI-talking sample}
\label{tab:real_outcome_dynamics_v31}
\caption*{\footnotesize \textit{Note.} This table reports forward operating and investment outcomes in the AI-talking annual sample. Coefficients are reported for the canonical grant-based PatentMismatch construct and alternative construct variants. Outcomes are measured one and two years after the disclosure year. Standard errors are shown in parentheses. *, **, and *** denote significance at the 10\%, 5\%, and 1\% levels, respectively.}
\resizebox{\textwidth}{!}{%
\begin{tabular}{lccccc}
\hline
Outcome & Canonical grant mismatch & Application mismatch & Low credibility only & Outcome mean & Observations \\
\hline
ROA t+1 & -0.0045 (0.0157) & 0.0003 (0.0156) & -0.0083 (0.0151) & -0.224 & 6865 \\
ROA t+2 & -0.0672** (0.0269) & -0.0526** (0.0267) & -0.0649*** (0.0248) & -0.223 & 4476 \\
Sales growth t+1 & 0.0231 (0.0502) & -0.0297 (0.0473) & -0.0065 (0.0430) & 0.240 & 6532 \\
Sales growth t+2 & -0.0240 (0.0550) & 0.0599 (0.0553) & 0.0550 (0.0460) & 0.211 & 4284 \\
CAPX/assets t+1 & -0.0003 (0.0009) & -0.0009 (0.0008) & -0.0006 (0.0008) & 0.024 & 6829 \\
CAPX/assets t+2 & 0.0012 (0.0010) & 0.0013 (0.0009) & 0.0014 (0.0010) & 0.023 & 4449 \\
R\&D/assets t+1 & 0.0097 (0.0070) & 0.0079 (0.0073) & 0.0087 (0.0064) & 0.149 & 4514 \\
R\&D/assets t+2 & 0.0178** (0.0070) & 0.0154** (0.0074) & 0.0166*** (0.0062) & 0.147 & 3044 \\
\hline
\end{tabular}
}
\end{table}


\subsection{Scrutiny without silence}

External scrutiny provides the first organizational-response test. Table~\ref{tab:comment_letter_cleanup_v32} studies within-firm changes around each firm's first SEC comment-letter year, restricting the main sample to active AI disclosers that mention AI in both $t-1$ and $t+1$. The key pattern is disclosure cleanup without silence. In the main sample, the AS ratio rises by $0.1796$ at $t+1$ ($p=0.007$), SpecShare falls by $0.1051$ at $t+2$ ($p=0.025$), and PatentMismatch declines by $0.1143$ at $t+1$ ($p=0.073$). AI Focus also rises after scrutiny rather than falling. The pattern is not explained by firms dropping AI language from their filings. Continuing AI disclosers rebalance toward more credible composition after scrutiny. The AI-specific comment subset is small, so the broader sample drives the estimates.

\begin{table}[!htbp]
\centering
\caption{Disclosure cleanup after SEC comment-letter scrutiny}
\label{tab:comment_letter_cleanup_v32}
\small
\resizebox{\textwidth}{!}{%
\begin{tabular}{lccccccl}
\hline
Outcome & t-2 vs t-1 & t vs t-1 & t+1 vs t-1 & t+2 vs t-1 & Pre mean & N firms \\
\hline
\multicolumn{7}{l}{\textit{Panel A. First SEC comment letter (active AI disclosers)}} \\
Speculative share & -0.1799*** (0.0519) & -0.0476 (0.0304) & -0.0621 (0.0407) & -0.1051** (0.0457) & 0.2717 & 70 \\
AS ratio & -0.1768*** (0.0644) & 0.0466 (0.0599) & 0.1796*** (0.0651) & 0.2213** (0.1049) & 0.7910 & 70 \\
PatentMismatch & -0.1935*** (0.0684) & -0.0429 (0.0516) & -0.1143* (0.0629) & -0.1207* (0.0699) & 0.4143 & 70 \\
AI Focus & -0.5977*** (0.0920) & 0.1993** (0.0794) & 0.5344*** (0.0964) & 0.5574*** (0.1409) & 1.4720 & 70 \\
\hline
\multicolumn{7}{l}{\textit{Panel B. First AI-related SEC comment letter (active AI disclosers)}} \\
Speculative share & -0.0097 (0.0750) & 0.0241 (0.0543) & 0.1360 (0.1034) & 0.0697 (0.1369) & 0.1471 & 11 \\
AS ratio & -0.1434 (0.1171) & 0.1883 (0.1213) & -0.0054 (0.1842) & 0.1767 (0.2122) & 0.8552 & 11 \\
PatentMismatch & -0.0909 (0.2113) & 0.0909 (0.0909) & 0.2727* (0.1408) & 0.1250 (0.2266) & 0.2727 & 11 \\
AI Focus & -0.4144** (0.1367) & 0.6666*** (0.1303) & 0.8162** (0.3076) & 1.1843** (0.3764) & 2.0795 & 11 \\
\hline
\end{tabular}
}
\begin{flushleft}
\footnotesize \textit{Note.} Each cell reports the mean within-firm change relative to the filing year immediately before the first comment-letter event ($t-1$) among active AI disclosers, defined as firms that mention AI in both $t-1$ and $t+1$. Standard errors are based on paired differences across firms. AI-related SEC comment uses the broad phrase match from Test 19 and should be read as a pilot because the event count is small. *, **, and *** denote significance at the 10\%, 5\%, and 1\% levels, respectively.
\end{flushleft}
\end{table}


\subsection{Regulatory salience and disclosure cleanup}

A second scrutiny test centers on the SEC's March 18, 2024 AI-washing enforcement actions (SEC, 2024). Table~\ref{tab:sec_enforcement_did_v32} fixes treatment using firms that were already exposed in 2022. Pre-exposed firms become less speculative and less likely to be flagged as mismatch after 2024. The post-2024 coefficient is $-0.0839$ for SpecShare, $0.2907$ for the AS ratio, and $-0.3121$ for PatentMismatch, all with $p<0.001$. The same pattern appears under alternative exposure definitions based on LowCredibility and ApplicationMismatch. The 2023 placebo estimates are already directionally similar for the composition outcomes, limiting a sharp enforcement interpretation. The post-2024 estimates are better read as regulatory salience or acceleration in disclosure cleanup. Substantively, disclosure composition becomes more credible after the SEC AI-washing salience event without a comparable collapse in AI Focus.

\begin{table}[!htbp]
\centering
\caption{SEC AI-washing enforcement salience and disclosure cleanup}
\label{tab:sec_enforcement_did_v32}
\small
\begin{tabular}{llccc}
\hline
\multicolumn{5}{l}{\textit{Panel A. Enforcement-salience DID using mismatch exposure in 2022}} \\
 & Outcome & Post-2024 DID & 2023 placebo & Mean / treated \\
\hline
 & Speculative share & -0.0839*** (0.0184) & -0.1224*** (0.0187) & 0.160 \\
 & AS ratio & 0.2907*** (0.0345) & 0.3314*** (0.0315) & 0.939 \\
 & PatentMismatch & -0.3121*** (0.0288) & -0.4305*** (0.0284) & 0.261 \\
 & AI Focus & -0.1044 (0.0653) & 0.0118 (0.0495) & 1.850 \\
\hline
\multicolumn{5}{l}{\textit{Panel B. Alternative exposure definitions on core disclosure outcomes}} \\
 & Treatment definition & SpecShare DID & AS ratio DID & PatentMismatch DID \\
\hline
 & LowCredibility in 2022 & -0.1215*** (0.0166) & 0.3368*** (0.0331) & -0.2606*** (0.0259) \\
 & ApplicationMismatch in 2022 & -0.1042*** (0.0185) & 0.3155*** (0.0344) & -0.2634*** (0.0290) \\
\hline
\end{tabular}
\begin{flushleft}
\footnotesize \textit{Note.} Treatment is fixed using firms that mention AI in 2022 and satisfy the exposure condition listed in the table. The post period is 2024--2025, after the SEC's March 18, 2024 AI-washing enforcement actions (SEC, 2024). Panel A reports the main PatentMismatch-exposure DID estimates and the 2023 placebo coefficient from the event-time specification. Panel B reports post-2024 DID coefficients for alternative exposure definitions on speculative share, AS ratio, and PatentMismatch, respectively. All specifications absorb firm and year fixed effects, include lagged controls, and cluster standard errors by firm. *, **, and *** denote significance at the 10\%, 5\%, and 1\% levels, respectively.
\end{flushleft}
\end{table}


\subsection{Managerial incentives and financing-window opportunism}

Managerial incentives show a related pattern. Table~\ref{tab:exec_incentive_mismatch_v32} links the AI-talking annual panel to lagged CEO observations from ExecuComp. In the post-ChatGPT sample, PatentMismatch is positively associated with lagged CEO equity-award share and CEO ownership, with coefficients of $0.1787$ ($p=0.041$) and $0.0170$ ($p=0.036$), respectively. In the big-firm sample, PatentMismatch is also positively associated with CEO ownership, with a coefficient of $0.1038$ ($p<0.001$). The estimates are not causal, but PatentMismatch is more prevalent when CEO incentives are more equity-oriented and managerial ownership is higher.

\begin{table}[!htbp]
\centering
\caption{Executive incentives and low-credibility AI disclosure}
\label{tab:exec_incentive_mismatch_v32}
\small
\resizebox{\textwidth}{!}{%
\begin{tabular}{llcccrr}
\hline
Panel & Outcome & Equity-award share & Ownership pct & Log CEO pay & Outcome mean & Obs. \\
\hline
\multicolumn{7}{l}{\textit{Panel A. ExecuComp-linked AI-talking sample}} \\
 & PatentMismatch & 0.0035 (0.0628) & 0.0052 (0.0101) & -0.0021 (0.0190) & 0.374 & 979 \\
 & LowCredibility & -0.0179 (0.0790) & -0.0008 (0.0084) & -0.0100 (0.0225) & 0.543 & 979 \\
 & AS ratio & 0.0896 (0.1189) & -0.0021 (0.0144) & 0.0677* (0.0350) & 0.785 & 979 \\
\multicolumn{7}{l}{\textit{Panel B. Big-firm ExecuComp sample}} \\
 & PatentMismatch & 0.0322 (0.0672) & 0.1038*** (0.0290) & 0.0003 (0.0198) & 0.375 & 771 \\
 & LowCredibility & 0.0002 (0.0912) & 0.0440 (0.0507) & -0.0108 (0.0242) & 0.541 & 771 \\
 & AS ratio & 0.0215 (0.1321) & -0.1113* (0.0674) & 0.0605 (0.0377) & 0.803 & 771 \\
\multicolumn{7}{l}{\textit{Panel C. Post-ChatGPT ExecuComp sample}} \\
 & PatentMismatch & 0.1787** (0.0874) & 0.0170** (0.0081) & 0.0528* (0.0321) & 0.442 & 407 \\
 & LowCredibility & 0.2036 (0.1247) & 0.0290*** (0.0065) & 0.0709 (0.0462) & 0.580 & 407 \\
 & AS ratio & -0.2022 (0.1255) & -0.0369*** (0.0069) & -0.0687 (0.0512) & 0.705 & 407 \\
\hline
\end{tabular}
}
\begin{flushleft}
\footnotesize \textit{Note.} The sample links the AI-talking annual panel to lagged CEO observations from ExecuComp, using the highest-TDC1 CEO row when multiple CEO-designated rows exist in a firm-year. Predictors are one-year lagged CEO equity-award share of total compensation, CEO ownership percentage excluding options, and log total CEO pay. Specifications absorb firm and year fixed effects, include the baseline annual control set, and cluster standard errors by firm. *, **, and *** denote significance at the 10\%, 5\%, and 1\% levels, respectively.
\end{flushleft}
\end{table}


Financing windows provide a sharper timing test. Table~\ref{tab:capital_raising_timing_v32} aligns disclosure dynamics around each firm's first large equity-issuance event and focuses on active AI issuers that continue discussing AI in the event year and the following year. Credibility deteriorates before the financing event and partly reverses afterward. Among all active AI issuers, LowCredibility rises by $0.2334$ and PatentMismatch rises by $0.1787$ into the issue year, then both decline in the following year. AI Focus, by contrast, continues to rise through and after the issue year. The non-big subsample shows the same timing shape. This timing is difficult to reconcile with silence or random variation alone. The evidence is more consistent with AI promotion around financing events followed by partial cleanup.

\begin{table}[!htbp]
\centering
\caption{Capital-raising timing and low-credibility AI disclosure}
\label{tab:capital_raising_timing_v32}
\small
\resizebox{\textwidth}{!}{%
\begin{tabular}{llcccrr}
\hline
Panel & Outcome & Event year vs. t-1 & t+1 vs. event year & t+2 vs. event year & Event-year mean & Obs. \\
\hline
\multicolumn{7}{l}{\textit{Panel A. First large equity issue among active AI issuers}} \\
 & LowCredibility & 0.2334*** (0.0294) & -0.1063*** (0.0233) & -0.1499*** (0.0297) & 0.4741 & 367 \\
 & ApplicationMismatch & 0.1816*** (0.0276) & -0.0790*** (0.0242) & -0.1008*** (0.0288) & 0.3597 & 367 \\
 & PatentMismatch & 0.1787*** (0.0269) & -0.0926*** (0.0242) & -0.1144*** (0.0285) & 0.3678 & 367 \\
 & Speculative share & 0.1187*** (0.0158) & -0.0177 (0.0118) & -0.0212 (0.0160) & 0.1813 & 367 \\
 & AI Focus & 1.0051*** (0.0545) & 0.3962*** (0.0410) & 0.4644*** (0.0592) & 1.6775 & 367 \\
\multicolumn{7}{l}{\textit{Panel B. Non-big first-issue AI issuers}} \\
 & LowCredibility & 0.2250*** (0.0344) & -0.0956*** (0.0276) & -0.1355*** (0.0360) & 0.4382 & 251 \\
 & ApplicationMismatch & 0.1875*** (0.0335) & -0.0797*** (0.0289) & -0.1076*** (0.0348) & 0.3705 & 251 \\
 & PatentMismatch & 0.1708*** (0.0323) & -0.0797*** (0.0289) & -0.1036*** (0.0342) & 0.3586 & 251 \\
 & Speculative share & 0.1231*** (0.0189) & -0.0171 (0.0144) & -0.0305 (0.0190) & 0.1832 & 251 \\
 & AI Focus & 1.0540*** (0.0663) & 0.3882*** (0.0507) & 0.4174*** (0.0766) & 1.6373 & 251 \\
\hline
\end{tabular}
}
\begin{flushleft}
\footnotesize \textit{Note.} Event year 0 is the disclosure year preceding a large equity-issuance window, where Issue $>5\%$ is defined as next-year CRSP shares-outstanding growth above 5\%. The event sample keeps each firm's first such issue year and focuses on active AI issuers that continue talking about AI in the event year and the following year. Positive values in the first comparison indicate deterioration or intensification into the issuance window; negative values in the second and third comparisons indicate partial unwind after the issue year. *, **, and *** denote significance at the 10\%, 5\%, and 1\% levels, respectively.
\end{flushleft}
\end{table}


\FloatBarrier

\subsection{Mixed and benchmark-sensitive market evidence}

The market-pricing evidence is more limited than the validation, scrutiny, and incentive evidence. Filing-date abnormal returns remain muted, so the data do not support a clean immediate pricing response at the annual 10-K filing date. Table~\ref{tab:factor_adjusted_alpha_v31} reports the benchmark-adjusted portfolio results. The strongest return result is the equal-weight three-month long-credible, short-mismatch portfolio, which delivers a FF5+momentum alpha of $0.976$ percentage points per month ($p=0.026$). The corresponding value-weight specification is statistically and economically close to zero. The return result is narrow and concentrated outside the largest firms.

Appendix results show that the short-run return relation depends on financing state. In those specifications, mismatch BHAR is materially less negative inside financing windows than outside them. Market pricing remains mixed and sensitive to benchmarks. The evidence is consistent with limited and state-dependent pricing of disclosure credibility, but it does not support a broad anomaly interpretation.

\begin{table}[!htbp]
\centering
\caption{Factor-adjusted calendar-time long-short alpha after AI filing signals}
\label{tab:factor_adjusted_alpha_v31}
\caption*{\footnotesize \textit{Note.} This table reports calendar-time long--short portfolio returns using the latest available annual filing signal for each stock-month. The long leg holds non-mismatch AI filers and the short leg holds mismatch AI filers. Panel A reports equal-weight portfolios and Panel B reports value-weight portfolios. Alpha columns report monthly intercepts from CAPM, Fama--French three-factor, Fama--French five-factor, and five-factor-plus-momentum regressions; p-values are shown in brackets.}
\resizebox{\textwidth}{!}{%
\begin{tabular}{lccccccc}
\hline
Horizon & Mean monthly & Annualized simple & CAPM alpha [p] & FF3 alpha [p] & FF5 alpha [p] & FF5+Mom alpha [p] & Months \\
\hline
\multicolumn{8}{l}{\textit{Panel A. Equal-weight long credible / short mismatch portfolios}} \\
1m & -0.508 & -6.102 & -0.218 [0.807] & 0.191 [0.829] & 0.183 [0.831] & 0.345 [0.675] & 98 \\
3m & 0.873 & 10.475 & 0.934 [0.036] & 1.043 [0.023] & 0.965 [0.032] & 0.976 [0.026] & 105 \\
6m & -0.033 & -0.393 & -0.015 [0.951] & -0.020 [0.941] & -0.003 [0.990] & -0.072 [0.803] & 106 \\
12m & -0.195 & -2.342 & -0.198 [0.246] & -0.170 [0.326] & -0.173 [0.333] & -0.176 [0.322] & 106 \\
\hline
\multicolumn{8}{l}{\textit{Panel B. Value-weight long credible / short mismatch portfolios}} \\
1m & 0.062 & 0.741 & 0.061 [0.943] & 0.122 [0.887] & 0.052 [0.950] & 0.129 [0.871] & 98 \\
3m & -0.138 & -1.658 & -0.007 [0.986] & 0.024 [0.954] & 0.053 [0.897] & 0.001 [0.997] & 105 \\
6m & 0.045 & 0.542 & -0.137 [0.666] & -0.089 [0.774] & -0.047 [0.885] & -0.032 [0.919] & 106 \\
12m & -0.185 & -2.218 & -0.364 [0.070] & -0.266 [0.167] & -0.198 [0.290] & -0.158 [0.410] & 106 \\
\hline
\end{tabular}
}
\end{table}


\FloatBarrier

The appendix reports fuller market checks, issue-window refinements, additional governance and intermediary tests, and measurement and sample-construction diagnostics.
